% Mitigation insertion for RE_SM_2611.tex.
% Intended location: immediately before \end{document}. This file is a fragment,
% not a replacement for the canonical Supporting Information. Insert with
% % Mitigation insertion for RE_SM_2611.tex.
% Intended location: immediately before \end{document}. This file is a fragment,
% not a replacement for the canonical Supporting Information. Insert with
% % Mitigation insertion for RE_SM_2611.tex.
% Intended location: immediately before \end{document}. This file is a fragment,
% not a replacement for the canonical Supporting Information. Insert with
% % Mitigation insertion for RE_SM_2611.tex.
% Intended location: immediately before \end{document}. This file is a fragment,
% not a replacement for the canonical Supporting Information. Insert with
% \input{insertions/RE_SM_mitigation_insertions_261114} and compile both documents.

\section{Weather-year and accepted-export sensitivity analyses}
\label{sec:sm_weather_export}

\subsection{ERA5 weather-year sensitivity, 2014--2023}
\label{sec:sm_weather_decade}

The main analysis uses 2023 hourly ERA5 forcing. To assess interannual weather sensitivity while holding the crop, optical, and economic parameters fixed, we repeated the WP3 simulations for each year from 2014 through 2023 at the three modeled sites. Leap years contain 8784 hourly samples; other years contain 8760. These are deterministic historical-reanalysis scenarios, not independent random draws, probabilities, confidence intervals, or field validation. The 2023 baseline outputs were retained separately.

\begin{table}[htbp]
  \centering
  \scriptsize
  \caption{Focal ST-OPV modeled annual crop-yield change and shade-nearest OPV--c-Si yield contrast across 2014--2023 ERA5 forcing. Means and standard deviations are across ten deterministic weather years; positive-year counts are descriptive.}
  \label{tab:weather_decade}
  \begin{tabular}{@{}lrrrrrr@{}}
    \toprule
    Crop & Mean $\Delta Y_{\rm OPV}$ & SD & Range & Positive years & Mean OPV--c-Si contrast & Positive contrast years \\
    & (\unit{\percent}) & (pp) & (\unit{\percent}) & & (pp) & \\
    \midrule
    Arabica coffee & $-8.22$ & 1.11 & $[-11.45,-6.38]$ & 0/10 & $-0.74$ & 0/10 \\
    Penja white pepper & $+1.85$ & 2.01 & $[-1.43,+7.60]$ & 7/10 & $+0.17$ & 8/10 \\
    Highland tea & $-3.54$ & 0.97 & $[-5.85,-1.86]$ & 0/10 & $-1.11$ & 0/10 \\
    \bottomrule
  \end{tabular}
\end{table}

For Penja white pepper, the shade-nearest OPV--c-Si annual yield contrast ranges from $-0.06$ to $+0.50$ percentage points, crossing zero. The positive-year counts apply only to this ten-year, fixed-parameter reanalysis set and should not be interpreted as probabilities of future outcomes. This sensitivity addresses weather-year variability; it does not sample uncertainty in crop parameters, soil properties, future climate, or model structure.

The focal design is approximately shade-matched to the c-Si reference, not energy-matched. Across 2014--2023, ST-OPV generation was \qtyrange{81.2}{82.8}{\percent} of generation for the shade-nearest c-Si design; the 2023 ratio was \qtyrange{81.0}{83.0}{\percent}. No tested equal-shade pair also matched annual electricity generation. Matching c-Si output would require reducing its ground-coverage ratio to approximately 30--35\%, thereby changing crop-light exposure. The comparison therefore matches crop-light conditions approximately, not both agricultural and electrical services.

\subsection{Accepted-export accounting sensitivity}
\label{sec:sm_export_acceptance}

The full-export baseline assumes acceptance of all otherwise-exportable electricity. In the sensitivity case, the hourly dispatch accepts \qty{50}{\percent} of otherwise-exportable annual energy and curtails the remainder; local consumption and the crop-specific selected battery configuration are retained. This annual acceptance fraction is not an hourly interconnection power limit, a grid-outage model, or evidence of a local purchase agreement. NPV is reported per hectare of installed array under the modeled cooperative-load allocation.

\begin{table}[htbp]
  \centering
  \scriptsize
  \caption{Modeled NPV for the focal ST-OPV scenario under full, \qty{50}{\percent}-accepted, and zero grid export. Each crop uses its selected battery configuration. NPV is per installed-array hectare; export acceptance is a stylized annual energy constraint.}
  \label{tab:export_acceptance}
  \begin{tabular}{@{}llrrr@{}}
    \toprule
    Crop & Selected battery & Full export & \qty{50}{\percent} accepted & No export \\
    & & \multicolumn{3}{c}{NPV (\unit{\USD\per\hectare})} \\
    \midrule
    Arabica coffee & Small LFP, \qty{2}{\hour} & 403,543.53 & 11,895.55 & $-379,752.43$ \\
    Penja white pepper & Small LFP, \qty{2}{\hour} & 368,372.49 & 43,412.03 & $-281,548.43$ \\
    Highland tea & Overnight withering BESS & 860,825.08 & 590,309.49 & 319,793.89 \\
    \bottomrule
  \end{tabular}
\end{table}

At \qty{50}{\percent} acceptance, rejected otherwise-exportable energy is curtailed. In these modeled cases, NPV declines from full export to partial acceptance and again to no export; coffee and pepper become unprofitable in the no-export cases, whereas tea remains positive. The results show that export access is decision-relevant under the assumptions tested, but do not establish that the local grid can accept the modeled energy or provide the assumed tariff. Site-specific interconnection and tariff data are needed before interpreting these values as project returns.
 and compile both documents.

\section{Weather-year and accepted-export sensitivity analyses}
\label{sec:sm_weather_export}

\subsection{ERA5 weather-year sensitivity, 2014--2023}
\label{sec:sm_weather_decade}

The main analysis uses 2023 hourly ERA5 forcing. To assess interannual weather sensitivity while holding the crop, optical, and economic parameters fixed, we repeated the WP3 simulations for each year from 2014 through 2023 at the three modeled sites. Leap years contain 8784 hourly samples; other years contain 8760. These are deterministic historical-reanalysis scenarios, not independent random draws, probabilities, confidence intervals, or field validation. The 2023 baseline outputs were retained separately.

\begin{table}[htbp]
  \centering
  \scriptsize
  \caption{Focal ST-OPV modeled annual crop-yield change and shade-nearest OPV--c-Si yield contrast across 2014--2023 ERA5 forcing. Means and standard deviations are across ten deterministic weather years; positive-year counts are descriptive.}
  \label{tab:weather_decade}
  \begin{tabular}{@{}lrrrrrr@{}}
    \toprule
    Crop & Mean $\Delta Y_{\rm OPV}$ & SD & Range & Positive years & Mean OPV--c-Si contrast & Positive contrast years \\
    & (\unit{\percent}) & (pp) & (\unit{\percent}) & & (pp) & \\
    \midrule
    Arabica coffee & $-8.22$ & 1.11 & $[-11.45,-6.38]$ & 0/10 & $-0.74$ & 0/10 \\
    Penja white pepper & $+1.85$ & 2.01 & $[-1.43,+7.60]$ & 7/10 & $+0.17$ & 8/10 \\
    Highland tea & $-3.54$ & 0.97 & $[-5.85,-1.86]$ & 0/10 & $-1.11$ & 0/10 \\
    \bottomrule
  \end{tabular}
\end{table}

For Penja white pepper, the shade-nearest OPV--c-Si annual yield contrast ranges from $-0.06$ to $+0.50$ percentage points, crossing zero. The positive-year counts apply only to this ten-year, fixed-parameter reanalysis set and should not be interpreted as probabilities of future outcomes. This sensitivity addresses weather-year variability; it does not sample uncertainty in crop parameters, soil properties, future climate, or model structure.

The focal design is approximately shade-matched to the c-Si reference, not energy-matched. Across 2014--2023, ST-OPV generation was \qtyrange{81.2}{82.8}{\percent} of generation for the shade-nearest c-Si design; the 2023 ratio was \qtyrange{81.0}{83.0}{\percent}. No tested equal-shade pair also matched annual electricity generation. Matching c-Si output would require reducing its ground-coverage ratio to approximately 30--35\%, thereby changing crop-light exposure. The comparison therefore matches crop-light conditions approximately, not both agricultural and electrical services.

\subsection{Accepted-export accounting sensitivity}
\label{sec:sm_export_acceptance}

The full-export baseline assumes acceptance of all otherwise-exportable electricity. In the sensitivity case, the hourly dispatch accepts \qty{50}{\percent} of otherwise-exportable annual energy and curtails the remainder; local consumption and the crop-specific selected battery configuration are retained. This annual acceptance fraction is not an hourly interconnection power limit, a grid-outage model, or evidence of a local purchase agreement. NPV is reported per hectare of installed array under the modeled cooperative-load allocation.

\begin{table}[htbp]
  \centering
  \scriptsize
  \caption{Modeled NPV for the focal ST-OPV scenario under full, \qty{50}{\percent}-accepted, and zero grid export. Each crop uses its selected battery configuration. NPV is per installed-array hectare; export acceptance is a stylized annual energy constraint.}
  \label{tab:export_acceptance}
  \begin{tabular}{@{}llrrr@{}}
    \toprule
    Crop & Selected battery & Full export & \qty{50}{\percent} accepted & No export \\
    & & \multicolumn{3}{c}{NPV (\unit{\USD\per\hectare})} \\
    \midrule
    Arabica coffee & Small LFP, \qty{2}{\hour} & 403,543.53 & 11,895.55 & $-379,752.43$ \\
    Penja white pepper & Small LFP, \qty{2}{\hour} & 368,372.49 & 43,412.03 & $-281,548.43$ \\
    Highland tea & Overnight withering BESS & 860,825.08 & 590,309.49 & 319,793.89 \\
    \bottomrule
  \end{tabular}
\end{table}

At \qty{50}{\percent} acceptance, rejected otherwise-exportable energy is curtailed. In these modeled cases, NPV declines from full export to partial acceptance and again to no export; coffee and pepper become unprofitable in the no-export cases, whereas tea remains positive. The results show that export access is decision-relevant under the assumptions tested, but do not establish that the local grid can accept the modeled energy or provide the assumed tariff. Site-specific interconnection and tariff data are needed before interpreting these values as project returns.
 and compile both documents.

\section{Weather-year and accepted-export sensitivity analyses}
\label{sec:sm_weather_export}

\subsection{ERA5 weather-year sensitivity, 2014--2023}
\label{sec:sm_weather_decade}

The main analysis uses 2023 hourly ERA5 forcing. To assess interannual weather sensitivity while holding the crop, optical, and economic parameters fixed, we repeated the WP3 simulations for each year from 2014 through 2023 at the three modeled sites. Leap years contain 8784 hourly samples; other years contain 8760. These are deterministic historical-reanalysis scenarios, not independent random draws, probabilities, confidence intervals, or field validation. The 2023 baseline outputs were retained separately.

\begin{table}[htbp]
  \centering
  \scriptsize
  \caption{Focal ST-OPV modeled annual crop-yield change and shade-nearest OPV--c-Si yield contrast across 2014--2023 ERA5 forcing. Means and standard deviations are across ten deterministic weather years; positive-year counts are descriptive.}
  \label{tab:weather_decade}
  \begin{tabular}{@{}lrrrrrr@{}}
    \toprule
    Crop & Mean $\Delta Y_{\rm OPV}$ & SD & Range & Positive years & Mean OPV--c-Si contrast & Positive contrast years \\
    & (\unit{\percent}) & (pp) & (\unit{\percent}) & & (pp) & \\
    \midrule
    Arabica coffee & $-8.22$ & 1.11 & $[-11.45,-6.38]$ & 0/10 & $-0.74$ & 0/10 \\
    Penja white pepper & $+1.85$ & 2.01 & $[-1.43,+7.60]$ & 7/10 & $+0.17$ & 8/10 \\
    Highland tea & $-3.54$ & 0.97 & $[-5.85,-1.86]$ & 0/10 & $-1.11$ & 0/10 \\
    \bottomrule
  \end{tabular}
\end{table}

For Penja white pepper, the shade-nearest OPV--c-Si annual yield contrast ranges from $-0.06$ to $+0.50$ percentage points, crossing zero. The positive-year counts apply only to this ten-year, fixed-parameter reanalysis set and should not be interpreted as probabilities of future outcomes. This sensitivity addresses weather-year variability; it does not sample uncertainty in crop parameters, soil properties, future climate, or model structure.

The focal design is approximately shade-matched to the c-Si reference, not energy-matched. Across 2014--2023, ST-OPV generation was \qtyrange{81.2}{82.8}{\percent} of generation for the shade-nearest c-Si design; the 2023 ratio was \qtyrange{81.0}{83.0}{\percent}. No tested equal-shade pair also matched annual electricity generation. Matching c-Si output would require reducing its ground-coverage ratio to approximately 30--35\%, thereby changing crop-light exposure. The comparison therefore matches crop-light conditions approximately, not both agricultural and electrical services.

\subsection{Accepted-export accounting sensitivity}
\label{sec:sm_export_acceptance}

The full-export baseline assumes acceptance of all otherwise-exportable electricity. In the sensitivity case, the hourly dispatch accepts \qty{50}{\percent} of otherwise-exportable annual energy and curtails the remainder; local consumption and the crop-specific selected battery configuration are retained. This annual acceptance fraction is not an hourly interconnection power limit, a grid-outage model, or evidence of a local purchase agreement. NPV is reported per hectare of installed array under the modeled cooperative-load allocation.

\begin{table}[htbp]
  \centering
  \scriptsize
  \caption{Modeled NPV for the focal ST-OPV scenario under full, \qty{50}{\percent}-accepted, and zero grid export. Each crop uses its selected battery configuration. NPV is per installed-array hectare; export acceptance is a stylized annual energy constraint.}
  \label{tab:export_acceptance}
  \begin{tabular}{@{}llrrr@{}}
    \toprule
    Crop & Selected battery & Full export & \qty{50}{\percent} accepted & No export \\
    & & \multicolumn{3}{c}{NPV (\unit{\USD\per\hectare})} \\
    \midrule
    Arabica coffee & Small LFP, \qty{2}{\hour} & 403,543.53 & 11,895.55 & $-379,752.43$ \\
    Penja white pepper & Small LFP, \qty{2}{\hour} & 368,372.49 & 43,412.03 & $-281,548.43$ \\
    Highland tea & Overnight withering BESS & 860,825.08 & 590,309.49 & 319,793.89 \\
    \bottomrule
  \end{tabular}
\end{table}

At \qty{50}{\percent} acceptance, rejected otherwise-exportable energy is curtailed. In these modeled cases, NPV declines from full export to partial acceptance and again to no export; coffee and pepper become unprofitable in the no-export cases, whereas tea remains positive. The results show that export access is decision-relevant under the assumptions tested, but do not establish that the local grid can accept the modeled energy or provide the assumed tariff. Site-specific interconnection and tariff data are needed before interpreting these values as project returns.
 and compile both documents.

\section{Weather-year and accepted-export sensitivity analyses}
\label{sec:sm_weather_export}

\subsection{ERA5 weather-year sensitivity, 2014--2023}
\label{sec:sm_weather_decade}

The main analysis uses 2023 hourly ERA5 forcing. To assess interannual weather sensitivity while holding the crop, optical, and economic parameters fixed, we repeated the WP3 simulations for each year from 2014 through 2023 at the three modeled sites. Leap years contain 8784 hourly samples; other years contain 8760. These are deterministic historical-reanalysis scenarios, not independent random draws, probabilities, confidence intervals, or field validation. The 2023 baseline outputs were retained separately.

\begin{table}[htbp]
  \centering
  \scriptsize
  \caption{Focal ST-OPV modeled annual crop-yield change and shade-nearest OPV--c-Si yield contrast across 2014--2023 ERA5 forcing. Means and standard deviations are across ten deterministic weather years; positive-year counts are descriptive.}
  \label{tab:weather_decade}
  \begin{tabular}{@{}lrrrrrr@{}}
    \toprule
    Crop & Mean $\Delta Y_{\rm OPV}$ & SD & Range & Positive years & Mean OPV--c-Si contrast & Positive contrast years \\
    & (\unit{\percent}) & (pp) & (\unit{\percent}) & & (pp) & \\
    \midrule
    Arabica coffee & $-8.22$ & 1.11 & $[-11.45,-6.38]$ & 0/10 & $-0.74$ & 0/10 \\
    Penja white pepper & $+1.85$ & 2.01 & $[-1.43,+7.60]$ & 7/10 & $+0.17$ & 8/10 \\
    Highland tea & $-3.54$ & 0.97 & $[-5.85,-1.86]$ & 0/10 & $-1.11$ & 0/10 \\
    \bottomrule
  \end{tabular}
\end{table}

For Penja white pepper, the shade-nearest OPV--c-Si annual yield contrast ranges from $-0.06$ to $+0.50$ percentage points, crossing zero. The positive-year counts apply only to this ten-year, fixed-parameter reanalysis set and should not be interpreted as probabilities of future outcomes. This sensitivity addresses weather-year variability; it does not sample uncertainty in crop parameters, soil properties, future climate, or model structure.

The focal design is approximately shade-matched to the c-Si reference, not energy-matched. Across 2014--2023, ST-OPV generation was \qtyrange{81.2}{82.8}{\percent} of generation for the shade-nearest c-Si design; the 2023 ratio was \qtyrange{81.0}{83.0}{\percent}. No tested equal-shade pair also matched annual electricity generation. Matching c-Si output would require reducing its ground-coverage ratio to approximately 30--35\%, thereby changing crop-light exposure. The comparison therefore matches crop-light conditions approximately, not both agricultural and electrical services.

\subsection{Accepted-export accounting sensitivity}
\label{sec:sm_export_acceptance}

The full-export baseline assumes acceptance of all otherwise-exportable electricity. In the sensitivity case, the hourly dispatch accepts \qty{50}{\percent} of otherwise-exportable annual energy and curtails the remainder; local consumption and the crop-specific selected battery configuration are retained. This annual acceptance fraction is not an hourly interconnection power limit, a grid-outage model, or evidence of a local purchase agreement. NPV is reported per hectare of installed array under the modeled cooperative-load allocation.

\begin{table}[htbp]
  \centering
  \scriptsize
  \caption{Modeled NPV for the focal ST-OPV scenario under full, \qty{50}{\percent}-accepted, and zero grid export. Each crop uses its selected battery configuration. NPV is per installed-array hectare; export acceptance is a stylized annual energy constraint.}
  \label{tab:export_acceptance}
  \begin{tabular}{@{}llrrr@{}}
    \toprule
    Crop & Selected battery & Full export & \qty{50}{\percent} accepted & No export \\
    & & \multicolumn{3}{c}{NPV (\unit{\USD\per\hectare})} \\
    \midrule
    Arabica coffee & Small LFP, \qty{2}{\hour} & 403,543.53 & 11,895.55 & $-379,752.43$ \\
    Penja white pepper & Small LFP, \qty{2}{\hour} & 368,372.49 & 43,412.03 & $-281,548.43$ \\
    Highland tea & Overnight withering BESS & 860,825.08 & 590,309.49 & 319,793.89 \\
    \bottomrule
  \end{tabular}
\end{table}

At \qty{50}{\percent} acceptance, rejected otherwise-exportable energy is curtailed. In these modeled cases, NPV declines from full export to partial acceptance and again to no export; coffee and pepper become unprofitable in the no-export cases, whereas tea remains positive. The results show that export access is decision-relevant under the assumptions tested, but do not establish that the local grid can accept the modeled energy or provide the assumed tariff. Site-specific interconnection and tariff data are needed before interpreting these values as project returns.
